本实验的主题是基本传递环节的分析。
在控制工程中，许多物理系统虽然来源不同，例如电气、机械或热学系统，但在线性化之后都可以用相同形式的传递函数描述。
因此，理解基本传递环节的时间响应和频率响应，是分析复杂控制系统的基础。

本实验研究 P、I、PT1 和 PT2 环节。
首先利用 MATLAB/Simulink 建立这些环节的模型，并观察它们在单位阶跃输入和正弦输入下的输出行为。
随后重点分析 PT2 环节，因为二阶系统已经能够表现出超调、振荡、阻尼衰减和共振等重要动态现象。

实验中使用阶跃响应评价系统在时域中的过渡过程，例如响应速度、稳态值、超调量和振荡衰减。
使用 Bode 图研究系统对不同频率正弦信号的幅值放大和相位偏移。
此外，通过零极点图将系统动态特性与复平面中的极点位置联系起来，从而更直观地判断稳定性、阻尼和固有频率的影响。

通过这些分析，可以建立从传递函数参数到系统实际动态行为之间的联系。
这为后续控制器设计、系统辨识和稳定性分析提供了必要的理论和仿真基础。
